\documentclass[english, parskip=full]{scrartcl}
\usepackage[utf8]{inputenc}  % Kodierung der Datei
\usepackage[english]{babel}  % Sprache auf Deutsch 
\usepackage{xcolor, colortbl}
\usepackage{graphicx, gensymb}
\usepackage{amsfonts, cancel, amssymb}
\usepackage{amsmath, amsthm, array, esint}
\usepackage{booktabs, multirow}  % schönere Tabellen
\usepackage{caption}  % Anpassung der Captions
\usepackage{scrlayer-scrpage}    % Manipulation des Seitenstils
\pagestyle{scrheadings}  % Stil für die Seite setzen
\automark{section}  % Abschnittsnamen als Seitenbeschriftung 
\setheadsepline{.5pt}    % Kopzeile durch Linie abtrennen
\usepackage[hidelinks]{hyperref}  % Links und weitere PDF-Features
\usepackage{typearea, tikz, pgffor, verbatim}
\usetikzlibrary{calc, arrows.meta,arrows, graphs, quotes}
\usepackage[cache=false, outputdir=build]{minted}
\usemintedstyle{vs}
\usepackage{placeins} % provides \FloatBarrier

\definecolor{bg}{rgb}{0.9,0.9,0.9}
\newcommand\numberthis{\addtocounter{equation}{1}\tag{\theequation}}
\usepackage{svg}
\usepackage{siunitx}
\usepackage{physics}
\usepackage{tensor}
\renewcommand{\bra}{}
\newcommand{\kom}{}
\newcommand{\brak}{}
\renewcommand{\braket}{}
\newcommand{\intx}{}
\newcommand{\intr}{}
\newcommand{\kla}{}
\newcommand{\klam}{}
\newcommand{\klammer}{}
\newcommand{\oper}{}
\newcommand{\tecxt}{}
\newcommand{\Bra}{}
\newcommand{\Ket}{}
\renewcommand{\i}{\text{i}}
\newcommand{\e}{\text{e}}
\newcommand*{\Fourier}{\textsc{Fourier}\ }
\newcommand*{\F}{\mathcal{F}}
\newcommand*{\sinc}{\text{sinc\,}}
\newcommand{\conv}{\mathop{\scalebox{3.5}{\raisebox{-0.38ex}{\makebox[0.5\width][c]{$\ast$}}}}\limits}

\newcommand*{\titel}{02 Coordinate systems}
\newcommand*{\autor}{Joachim Schwardt}

\hypersetup{pdfauthor={\autor}, pdftitle={\titel}}

%\titlehead{titlehead}
\subject{General Relativity}
\title{\titel}
\author{\autor}
\date{\today}

\begin{document}

%\begin{titlepage}
\maketitle  % Titel setzen
%\tableofcontents  % Inhaltsverzeichnis setzen
%\end{titlepage}


\section*{Task 4: Non-orthogonal coordinates}
Given is the vector $\vec{r}=\binom{x^1+x^2}{x^2}$.
We can compute $\vec{g}_a$ by using their definition
\begin{align}
\vec{g}_1 &= \partial_1\vec{r} = \binom{1}{0},&& &\vec{g}_2 &=\partial_2\vec{r} = \binom{1}{1} \label{eqn:g unten a} 
\end{align}
We may use the orthonormality condition $\vec{g}_a\cdot\vec{g}^b=\delta_a^b$. 
Let $\vec{g}^{1}=\binom{\alpha_{1}}{\beta_{1}}$ and $\vec{g}^{2}=\binom{\alpha_{2}}{\beta_{2}}$.
Then we find
\begin{align}
1 &\overset{!}{=} \vec{g}_1\cdot\vec{g}^1 = \alpha_1, &&& 0&\overset{!}{=}\vec{g}_2\cdot\vec{g}^1 = \alpha_1+\beta_1\ \ \Rightarrow\ \ \beta_1 = -1\\
0 &\overset{!}{=} \vec{g}_1\cdot\vec{g}^2 = \alpha_2, &&& 1&\overset{!}{=}\vec{g}_2\cdot\vec{g}^2 = \alpha_2+\beta_2\ \ \Rightarrow\ \ \beta_2 = 1
\end{align}
These results can be summarized as
\begin{align}
\vec{g}^1 &= g^{1b}\vec{g}_b = \binom{1}{-1},&& &\vec{g}^2 &= g^{2b}\vec{g}_b = \binom{0}{1}\label{eqn:g oben a}.
\end{align}
The vectors are visualized in \autoref{figure:g oben unten a}.
\begin{figure}[h!]
\centering
\begin{tikzpicture}[scale=2]
\draw[->, >=latex] (0,-1) -- (0,2) node[right]{$x^2$};
\draw[->, >=latex] (-2,0) -- (2,0) node[below]{$x^1$};
\draw[->, >=latex, blue, ultra thick] (0,0) -- (1,0) node[above]{$\vec{g}_1$};
\draw[->, >=latex, blue, ultra thick] (0,0) -- (1,1) node[above]{$\vec{g}_2$};
\draw[->, >=latex, red, ultra thick] (0,0) -- (1,-1) node[right]{$\vec{g}^1$};
\draw[->, >=latex, red, ultra thick] (0,0) -- (0,1) node[left]{$\vec{g}^2$};
\end{tikzpicture}
\caption{Visualization of $\vec{g}_a$ and $\vec{g}^a$.}
\label{figure:g oben unten a}
\end{figure}
The metric tensor and its inverse are given by
\begin{align}
g_{ab} &= \begin{pmatrix}
1&1 \\ 1&2
\end{pmatrix}&\Rightarrow&& g^{ab} &= \begin{pmatrix}
2&-1 \\ -1&1
\end{pmatrix} \label{eqn:Metrik}
\end{align}
The kinetic energy as a function of $\dot{x}^a$ is given by
\begin{align}
T &= \frac{m}{2}g_{ab}\dot{x}^a\dot{x}^b = \frac{m}{2}\kla\left[ {(\dot{x}^1)^2 + 2\dot{x}^1\dot{x}^2 + 2(\dot{x}^2)^2} \right] \label{eqn:Kinetische Energie von xdot oben a}
\end{align}
The gradient of some function $f(x^1,x^2)$ can be written as
\begin{align}
\nabla f &= \vec{g}^a\partial_a f = \binom{1}{-1}\partial_1f+\binom{0}{1}\partial_2f = \binom{\partial_1f}{-\partial_1f+\partial_2f} \label{eqn:Gradient}
\end{align}
We can compute the coefficients for $\vec{A}=\binom{1}{0}$ and $\vec{B}=\binom{0}{1}$ using
\begin{align}
A_1 &= \vec{A}\cdot\vec{g}_1 = 1 & A_2 &= \vec{A}\cdot\vec{g}_2 = 1,\\
A^1 &= \vec{A}\cdot\vec{g}^1 = 1 & A^2 &= \vec{A}\cdot\vec{g}^2 = 0,\\
B_1 &= \vec{B}\cdot\vec{g}_1 = 0 & B_2 &= \vec{B}\cdot\vec{g}_2 = 1,\\
B^1 &= \vec{B}\cdot\vec{g}^1 = -1 & B^2 &= \vec{B}\cdot\vec{g}^2 = 1.
\end{align}

\section*{Task 5: Cylindrical coordinates}
Given is the vector $\vec{r}=\begin{pmatrix}
x^1\cos x^2 \\ x^1\sin x^2 \\ x^3
\end{pmatrix}$ where $x^1=\rho$, $x^2=\phi$ and $x^3=z$.
We can compute $\vec{g}_a$ form the definition
\begin{align}
\vec{g}_1 &= \partial_1\vec{r} = \begin{pmatrix}
\cos x^2 \\ \sin x^2 \\ 0
\end{pmatrix} &\vec{g}_2 &= \partial_2\vec{r} = \begin{pmatrix}
-x^1\sin x^2 \\ x^1\cos x^2 \\ 0
\end{pmatrix}&\vec{g}_3 &= \partial_3\vec{r}=\begin{pmatrix}
0 \\ 0 \\ 1
\end{pmatrix} \label{eqn:zylinder:g unten a}
\end{align}
The metric tensor and its inverse are
\begin{align}
g_{ab} &= \begin{pmatrix}
1&0&0 \\ 0&(x^1)^2&0 \\ 0&0&1
\end{pmatrix} &\Rightarrow&& g^{ab} &= \begin{pmatrix}
1&0&0 \\ 0&\frac{1}{(x^1)^2}&0 \\ 0&0&1
\end{pmatrix}. \label{eqn:zylinder:Metrik}
\end{align}
We may use this to derive the $\vec{g}^a$ components
\begin{align}
\vec{g}^1 &= g^{1b}\vec{g}_b = \begin{pmatrix}
\cos x^2 \\ \sin x^2 \\ 0
\end{pmatrix} &\vec{g}^2 &= g^{2b}\vec{g}_b = \begin{pmatrix}
-\frac{1}{x^1}\sin x^2 \\ \frac{1}{x^1}\cos x^2 \\ 0
\end{pmatrix} &\vec{g}^3 &= g^{3b}\vec{g}_b = \begin{pmatrix}
0 \\ 0 \\ 1
\end{pmatrix}. \label{eqn:zylinder:g oben a}
\end{align}
The gradient of some function $f(x^1,x^2,x^3)$ may be written as
\begin{align}
\nabla f &= \vec{g}^a\partial_a f = \begin{pmatrix}
\cos x^2\partial_1f-\sin x^2\frac{\partial_2f}{x^1}\\ \cos x^2\partial_1f + \cos x^2\frac{\partial_2f}{x^1} \\ \partial_3f
\end{pmatrix}. \label{eqn:zylinder:Gradient}
\end{align}
The kinetic energy as a function of $\dot{x}^a$ is given by
\begin{align}
T &= \frac{m}{2}g_{ab}\dot{x}^a\dot{x}^b = \frac{m}{2}\kla\left[ {(\dot{x}^1)^2 + (x^1)^2(\dot{x}^2)^2 + (\dot{x}^3)^2} \right] = \\
&= \frac{m}{2} \klam\left[ \dot{\rho}^2 + \rho^2\dot{\phi}^2 + \dot{z}^2 \right]• \label{eqn:zylinder:Kinetische Energie}
\end{align}
The Christoffel symbol $\Gamma_{22}^1$ can be computed using the basis vectors
\begin{align}
\Gamma_{22}^1 &= \vec{g}^1\partial_2\vec{g}_2 = \begin{pmatrix}
\cos x^2 \\ \sin x^2 \\ 0
\end{pmatrix} \cdot \begin{pmatrix}
-x^1\cos x^2 \\ -x^1\sin x^2 \\ 0
\end{pmatrix} = -x^1
\end{align}
or by using the metric tensor explicitly
\begin{align}
\Gamma_{22}^1 &= \frac{1}{2}g^{d1}\kla\left( {g_{d2,2}+g_{2d,2}-g_{22,d}} \right) = -\frac{1}{2}g^{11}\partial_1(x^1)^2 = -x^1.
\end{align}
Note that one could also use the slightly simpler form $\Gamma^c_{ab} = g^{cd}\kla\left( g_{ad,b} - \frac{1}{2}g_{ab,d} \right)$, which only requires 2 terms (but makes the symmetry $a\leftrightarrow b$ less evident).
\\
A completely different way of getting at the Christoffel symbols is by first considering the Euler-Lagrange equations.
For this derivation we will use $(\rho,\phi,z)$ instead of $(x^1,x^2,x^3)$.
In our case we have $L\equiv T$ and therefore
\begin{align}
0 &= \frac{\partial L}{\partial \rho} - \frac{\tecxt\text{d}}{\tecxt\text{d}t} \frac{\partial L}{\partial \dot{\rho}} \propto \rho\dot{\phi}^2 - \ddot{\rho},\\
0 &= \frac{\partial L}{\partial \phi} - \frac{\tecxt\text{d}}{\tecxt\text{d}t} \frac{\partial L}{\partial \dot{\phi}} \propto 2\rho\dot{\rho}\dot{\phi} + \rho^2\ddot{\phi},\\
0 &= \frac{\partial L}{\partial z} - \frac{\tecxt\text{d}}{\tecxt\text{d}t} \frac{\partial L}{\partial \dot{z}} \propto \ddot{z}.
\end{align}
We know that $\ddot{x}^\alpha + \Gamma^\alpha_{\mu\nu} \dot{x}^\mu \dot{x}^\nu = 0$ and may thus read off the only non-zero Christoffel symbols as
\begin{align}
\Gamma^\rho_{\phi\phi} &= -\rho,\\
\Gamma^\phi_{\rho\phi} &= \frac{1}{\rho}.
\end{align}


%\begin{thebibliography}{9}
%\bibitem{example} description, \url{https://www.exmapleurl}, day.\,month~2021.
%\end{thebibliography}

\end{document}